\chapter{State of the Art and Theory}\label{chap:2}
{\itshape This chapter reviews the mathematical foundations of classical and quantum clustering. It defines the coordinates of our 5D color-spatial manifold, derives standard video combinatorics constraints, describes the CUDA parallel SIMT execution model, outlines the physics boundaries of physical NISQ processors, and presents the mathematical derivation of the quantum SWAP test phase simulator.}

\section{Classical K-Means Clustering and Spatial Weighting}
K-Means is a popular clustering algorithm that groups $N$ data points into $K$ uniform clusters. The algorithm works by iteratively minimizing WCSS:
\begin{equation}
\label{eq:wcss}
J = \sum_{i=1}^{K} \sum_{\mathbf{x} \in S_i} \|\mathbf{x} - \boldsymbol{\mu}_i\|^2
\end{equation}
where $S_i$ is the set of points in the $i$-th cluster, and $\boldsymbol{\mu}_i$ is the centroid of $S_i$. Each iteration has two main steps:
\begin{enumerate}
    \item \textbf{Assignment Step:} Every point $\mathbf{x}$ is assigned to its nearest centroid based on Euclidean distance: $S_i^{(t)} = \{ \mathbf{x} : \|\mathbf{x} - \boldsymbol{\mu}_i^{(t)}\|^2 \le \|\mathbf{x} - \boldsymbol{\mu}_j^{(t)}\|^2 \ \forall j \}$.
    \item \textbf{Update Step:} New centroids are calculated by averaging all points in each cluster: $\boldsymbol{\mu}_i^{(t+1)} = \frac{1}{|S_i^{(t)}|} \sum_{\mathbf{x} \in S_i^{(t)}} \mathbf{x}$.
\end{enumerate}
These steps repeat until the centroids stop changing or the algorithm reaches a maximum iteration limit ($I_{\text{max}}$).

\subsection{The 5D Color-Spatial Manifold}
Traditional clustering operates solely in 3D color spaces (such as RGB, HSV, or CIELAB). However, grouping pixels only by color causes spatial fragmentation, where pixels of similar color are grouped together regardless of their physical location in the image. To ensure spatial consistency, the feature space is extended to a 5D manifold:
\begin{equation}
\mathbf{x} = [B, G, R, x, y]^T
\end{equation}
where $B, G, R$ represent color channels, and $x, y$ are the pixel coordinates.

However, color values and spatial coordinates have completely different scales. In high-resolution images, the pixel coordinates would dominate the distance formula, ignoring color differences. To balance this, a spatial weight parameter ($\lambda_{\text{spatial}}$) is added: $\|\mathbf{x}_n - \boldsymbol{\mu}_i\|^2 = (B_n - \boldsymbol{\mu}_{i,B})^2 + (G_n - \boldsymbol{\mu}_{i,G})^2 + (R_n - \boldsymbol{\mu}_{i,R})^2 + \lambda_{\text{spatial}} \left[ (x_n - \boldsymbol{\mu}_{i,x})^2 + (y_n - \boldsymbol{\mu}_{i,y})^2 \right]$. A high $\lambda_{\text{spatial}}$ creates very compact, regular shapes, while a low value lets the engine prioritize color boundaries.

\subsection{Real-World Vision Applications and Edge Computing Constraints}
Real-time 5D clustering is useful for computer vision applications like lane tracking in self-driving cars, obstacle detection, or surgical assistance overlays. In automated vision loops, standard display interfaces refresh at $60\text{ FPS}$ ($16.67\text{ ms}$ processing budget), while high-speed industrial sensors or virtual reality setups require at least $120\text{ FPS}$ ($8.33\text{ ms}$ processing budget) to operate without control-loop latency.

Running big deep learning models on embedded devices like microcontrollers consumes too much power and has excessive latency. Lightweight, GPU-accelerated 5D clustering works as an efficient preprocessing step. It groups pixels quickly before sending compressed data to heavier algorithms, which satisfies the power and speed constraints of edge devices.

\subsection{Mathematical Clustering Quality Metrics}
Three standard metrics are used to evaluate the clustering results: WCSS, the Davies-Bouldin Index, and the Silhouette Coefficient.

\subsubsection{Within-Cluster Sum of Squares (WCSS)}
As defined in Equation~\ref{eq:wcss}, WCSS is the active objective function minimized by the classical K-Means algorithm during execution. While WCSS measures how compact a cluster is, it naturally decreases as more clusters $K$ are added. In this framework, WCSS is used as a core metric to evaluate how well both engines compress the feature space. However, because WCSS does not measure how well-separated different clusters are from each other, the Davies-Bouldin Index and the Silhouette Coefficient are also implemented as background telemetry diagnostics to compare the overall segmentation quality between the classical and quantum-inspired engines.

\subsubsection{The Davies-Bouldin Index}
The Davies-Bouldin Index ($\mathcal{D}\mathcal{B}$) evaluates clustering by checking the ratio of cluster size to the distance between cluster centers. We define cluster dispersion as $s_i = \left( \frac{1}{|S_i|} \sum_{\mathbf{x} \in S_i} \|\mathbf{x} - \boldsymbol{\mu}_i\|^2 \right)^{1/2}$ and centroid distance as $d(\boldsymbol{\mu}_i, \boldsymbol{\mu}_j) = \|\boldsymbol{\mu}_i - \boldsymbol{\mu}_j\|$. The DB index is the average similarity ratio across clusters: $\mathcal{D}\mathcal{B} = \frac{1}{K} \sum_{i=1}^{K} \max_{j \neq i} \left( \frac{s_i + s_j}{\|\boldsymbol{\mu}_i - \boldsymbol{\mu}_j\|} \right)$. A lower score means better, more separated clusters.

\subsubsection{The Silhouette Coefficient and Monte Carlo Approximation}
The Silhouette Coefficient $s(\mathbf{x}) \in [-1, 1]$ evaluates how well each point fits its cluster. For a point $\mathbf{x}$, its average distance to points in the same cluster ($a(\mathbf{x}) = \frac{1}{|S_A| - 1} \sum_{\mathbf{y} \in S_A, \mathbf{y} \neq \mathbf{x}} d(\mathbf{x}, \mathbf{y})$) and its distance to the nearest neighboring cluster ($b(\mathbf{x}) = \min_{C \neq A} \left( \frac{1}{|S_C|} \sum_{\mathbf{z} \in S_C} d(\mathbf{x}, \mathbf{z}) \right)$) are calculated, yielding the score: $s(\mathbf{x}) = \frac{b(\mathbf{x}) - a(\mathbf{x})}{\max(a(\mathbf{x}), b(\mathbf{x}))}$. The global score is the average across all points.

Calculating this exactly for a standard VGA video frame of $640 \times 480$ pixels ($N = 307,200$) is computationally impossible in real-time. Because the Silhouette score requires finding the pairwise distance between every single pixel, the total number of operations scales quadratically as $\frac{N(N-1)}{2}$, yielding $\frac{307,200 \times 307,199}{2} \approx 47.19$ billion distance calculations per frame. Executing 47 billion calculations in a real-time $120\text{ FPS}$ pipeline (with a strict $8.33\text{ ms}$ budget) is impossible. To bypass this bottleneck, a Monte Carlo approximation is used by sampling a random, uniform subset of $M=2,000$ pixels. This reduces the complexity to $\mathcal{O}(M^2)$ pairwise checks within the subset, which runs in less than a millisecond on a background thread, while remaining highly accurate.

\section{Parallel Computing and GPU Architectures via CUDA}

\subsection{Computational Complexity Constraints of Sequential Architectures}
Running K-Means sequentially on a CPU is too slow for real-time video. For a VGA frame with $N=307,200$ pixels, $K=8$ clusters, and $I=15$ iterations, the CPU must perform over $36$ million calculations per frame. To maintain a frame rate of $120\text{ FPS}$, the CPU has to do this in under $8.33\text{ ms}$, which is very difficult due to memory latency and single-core bottlenecks.

\subsection{The CUDA Execution Paradigm and Thread Hierarchy}
To solve this, the calculations are offloaded to the GPU using CUDA. CUDA uses the SIMT model, allowing thousands of threads to process different pixels at the same time. The GPU organizes these threads into a hierarchy:
\begin{itemize}
    \item \textbf{Grids and Blocks:} Launching a kernel initializes a \textit{Grid}, divided into \textit{Thread Blocks}. Threads in the same block run on the same Streaming Multiprocessor (SM), allowing them to share low-latency local cache and synchronize using hardware barriers.
    \item \textbf{Warps:} At the hardware level, threads execute in groups of $32$ called \textit{Warps}. If threads in a warp take different execution paths (like an \texttt{if-else} branch), they serialize. This warp divergence reduces GPU utilization, so keeping thread execution paths uniform is a key goal.
\end{itemize}

\subsection{GPU Memory Hierarchy and Optimization Strategies}
Efficiently utilizing the GPU requires managing its multi-tiered memory to avoid execution stalls:
\begin{itemize}
    \item \textbf{Global Memory:} This is the main VRAM. It is large but slow ($400$ to $800$ cycles). To speed it up, memory accesses are coalesced so threads read contiguous memory blocks in a single transaction.
    \item \textbf{Shared Memory (\texttt{\_\_shared\_\_}):} This is an on-chip, low-latency cache layer on the SM ($1$ to $2$ cycles). By caching the 5D cluster centroids in shared memory, repeatedly reading them from slow VRAM is avoided.
    \item \textbf{Registers:} These are private, zero-latency storage spaces for each thread. However, using too many registers per thread limits block concurrency on the SM, reducing overall GPU occupancy.
\end{itemize}

\subsection{Parallel Strided Downsampling Theory}
To keep performance stable across different GPUs, a strided downsampling preprocessor is implemented directly in the CUDA kernel. Instead of using CPU resizing, downsampling is performed as the image is loaded onto the GPU.

For a stride $S \in [1, 16] \subset \mathbb{Z}^+$, a naive implementation would launch a thread for every pixel and skip inactive ones using a modulo check: $\text{Active}(idx) = (idx \pmod{S} == 0)$. This causes warp divergence and slow memory reads. Instead, the GPU grid is scaled to the exact downsampled size $W_{\text{out}} = \lfloor W/S \rfloor$ and $H_{\text{out}} = \lfloor H/S \rfloor$. Each thread $(x_{\text{out}}, y_{\text{out}})$ maps directly to the input array at $x_{\text{in}} = x_{\text{out}} \cdot S$ and $y_{\text{in}} = y_{\text{out}} \cdot S$. This ensures all threads are active and eliminates divergence.

\section{Quantum Machine Learning and NISQ Limitations}

\subsection{The NISQ Era and Physical Hardware Constraints}
Quantum computing has interesting theoretical benefits, but current hardware is in the NISQ era. These processors have few qubits and suffer from dephasing ($T_2$ coherence) and relaxation ($T_1$ thermal decay), where qubits quickly lose their states within microseconds. They also have high gate error rates, particularly for two-qubit operations ($10^{-3}$ to $10^{-2}$). Deep quantum circuits accumulate too much noise to be useful.

\subsection{Hybrid Classical-Quantum Algorithms}
To work around these limits, hybrid classical-quantum algorithms (like VQE or QAOA) offload part of the work to a classical CPU or GPU. For clustering, the QPU calculates distances, while the classical processor updates the centroids. This hybrid loop allows for the utilization of quantum representations without exceeding the physical hardware limits.

\subsection{The Real-Time API and Simulation Latency Bottleneck}
However, using a real QPU or a python simulator in a live video pipeline is impossible due to massive bottlenecks:
\begin{itemize}
    \item \textbf{State Preparation Overhead:} Encoding pixel features into qubits takes too many gates, which exceeds the real-time frame window.
    \item \textbf{API and Network Latency:} Sending circuits to a remote QPU takes hundreds of milliseconds or seconds, which breaks real-time frame pacing.
    \item \textbf{Quantum Measurement Sampling Delay:} We have to repeat quantum measurements thousands of times (shots) to get stable probabilities.
\end{itemize}
Since real QPUs and simulators are too slow, quantum behaviors must be simulated deterministically using parallelized GPU kernels.

\section{Quantum-Inspired Computing and the SWAP Test}

\subsection{Quantum Overlap and the SWAP Test Circuit}
To get the mathematical benefits of quantum computing without the physical hardware latency, the quantum SWAP test is simulated directly on the GPU. In quantum information theory, the SWAP test measures the overlap (similarity) between two states $\lvert\psi\rangle$ and $\lvert\phi\rangle$ using three registers: an ancilla qubit initialized to $\lvert 0 \rangle_{\text{anc}}$, and two target registers holding the inputs. A Hadamard gate is applied to the ancilla, a Controlled-SWAP (CSWAP) gate swaps the targets if the ancilla is in state $\lvert1\rangle$, and a second Hadamard gate induces interference.

\subsection{Step-by-Step Mathematical Derivation}
The registers evolve as follows: the initial state is $\lvert\Psi_0\rangle = \lvert 0 \rangle_{\text{anc}}\lvert\psi\rangle\lvert\phi\rangle$. The first Hadamard gate yields the superposition $\lvert\Psi_1\rangle = \frac{1}{\sqrt{2}}\lvert0\rangle_{\text{anc}}\lvert\psi\rangle\lvert\phi\rangle + \frac{1}{\sqrt{2}}\lvert1\rangle_{\text{anc}}\lvert\psi\rangle\lvert\phi\rangle$. The CSWAP gate swaps the target registers if the ancilla is $\lvert1\rangle$, giving $\lvert\Psi_2\rangle = \frac{1}{\sqrt{2}}\lvert0\rangle_{\text{anc}}\lvert\psi\rangle\lvert\phi\rangle + \frac{1}{\sqrt{2}}\lvert1\rangle_{\text{anc}}\lvert\phi\rangle\lvert\psi\rangle$. A second Hadamard gate results in:
\begin{equation}
\label{eq:psi_3}
\lvert\Psi_3\rangle = \frac{1}{2}\lvert0\rangle_{\text{anc}}(\lvert\psi\rangle\lvert\phi\rangle + \lvert\phi\rangle\lvert\psi\rangle) + \frac{1}{2}\lvert1\rangle_{\text{anc}}(\lvert\psi\rangle\lvert\phi\rangle - \lvert\phi\rangle\lvert\psi\rangle)
\end{equation}

According to Born's Rule, the probability $P(\lvert 0 \rangle)$ of measuring the ancilla in the ground state is:
\begin{align}
P(\lvert0\rangle) &= \left\| \langle 0 \rvert \Psi_3 \rangle \right\|^2 = \left\| \frac{1}{2}(\lvert\psi\rangle\lvert\phi\rangle + \lvert\phi\rangle\lvert\psi\rangle) \right\|^2 \\
&= \frac{1}{4} \left[ \langle\psi\vert\psi\rangle\langle\phi\vert\phi\rangle + 2\lvert\langle\psi\vert\phi\rangle\rvert^2 + \langle\phi\vert\phi\rangle\langle\psi\vert\psi\rangle \right] = \frac{1}{2} + \frac{1}{2}\lvert\langle\psi\vert\phi\rangle\rvert^2
\end{align}
This proves that the ground state probability is directly related to the state overlap $\lvert\langle\psi\vert\phi\rangle\rvert^2$.

\subsection{Hilbert Phase Space Mapping and Implemented Metrics}
This state overlap is implemented directly in CUDA C++ without virtual quantum gates. The normalized 5D pixel features are mapped to phase states: $\theta_d = (p_d - c_d) \cdot \lambda_{\text{scale}}$, where $p_d$ is the pixel, $c_d$ is the centroid, and $\lambda_{\text{scale}}$ is a constant. 

The 1D overlap is $\text{Overlap}_d = \cos^2(\theta_d)$. The global similarity is the product of these overlaps: $\text{Fidelity}(\mathbf{p}, \mathbf{c}) = \prod_{d=1}^{D} \cos^2\left( (p_d - c_d) \cdot \lambda_{\text{scale}} \right)$. Since K-Means needs a minimizing distance, we subtract this similarity from one:
\begin{equation}
D_{\text{Quantum}}(\mathbf{p},\mathbf{c}) = 1.0 - \prod_{d=1}^{D} \cos^2\left( (p_d - c_d) \cdot \lambda_{\text{scale}} \right)
\end{equation}
The scaling coefficient is locked to $\lambda_{\text{scale}} = \text{SCALE\_FACTOR} \cdot \text{PHASE\_OFFSET} = \frac{\pi}{2} \cdot 0.5 = \frac{\pi}{4}$. This bounds the angle between $[0, \frac{\pi}{4}]$, allowing the GPU to calculate the cosine quickly using hardware-accelerated instructions (\texttt{\_\_cosf}).

\subsection{Scientific Justification of Phase Boundaries}
Using this non-linear trigonometric distance instead of the standard Euclidean norm has a real advantage for video segmentation. Standard Euclidean distance changes linearly. In image regions with gradual shadows or noise, this linear change can cause pixels to bleed across boundaries.

In contrast, the quantum-inspired metric changes along a squared cosine curve. As the distance between a pixel and a centroid grows, the derivative of $\cos^2(\theta)$ drops quickly. This creates a steep distance slope at the cluster edges, acting like a natural noise filter. Even in dark or noisy areas, these steep boundaries keep the cluster centers stable, resulting in much cleaner and sharper object segmentation. This mathematical foundation directly informs the structural design and class interactions of the software architecture detailed in the next chapter.
